\Page{12}{\textenglish{03 / CAMPAIGNS}}{ארבע תמונות, שלוש דרכי הפקה}{43}{49}

{\fontsize{12.6}{18}\selectfont

נניח שמותג צריך ארבע תמונות מאושרות: אותו בגד, אותה דוגמנית, אותו חודש ואותם ערוצים. בהפקת צילום, הצוות מביא את הדוגמנית, הבגד והמקום אל המצלמה. בתהליך משולב הוא יכול לצלם את הדוגמנית בבגד ולהכין לה סביבה אחרת בפוסט. בתהליך מחולל, חלק גדול יותר מהפריים נבנה אחרי הכנת חומרי ייחוס. בריף אחד יכול להוביל לשלוש רשימות עבודה שונות, אפילו אם הארבע שיגיעו למגזין נראות כאילו צולמו באותו אחר צהריים.

\par}

\Photo{hybrid-editorial.png}{77mm}{סט אופנה מדומיין להמחשת קומפוזיציה. \AIArt}

{\fontsize{12.2}{17.2}\selectfont

\StoryColumns{בקיץ 2024, בקמפיין \textenglish{Sunset Dream} של \textenglish{Mango Teen}, המצלמה שימשה תחילה לצילום פריטי הקולקציה. לפי החברה, הצילומים שימשו לאימון מודל שהציב את הבגדים על דוגמנית בדימוי מחולל. מכאן באה עבודה מוכרת של בחירה, ריטוש ועריכה. אנשי העיצוב, האמנות והסטיילינג עבדו לצד אנשי הנתונים ואימון המודל ובשיתוף סטודיו הצילום. האתגר שתיארה החברה היה לשמור על מאפייני הבגד והדוגמנית גם כשהפריים עצמו נוצר בכלי אחר. \SourceRef{7}}{להפקה כזאת יש סדר עבודה משלה. אפשר להתחיל מסקיצה של קומפוזיציה או מתיאור מילולי, לארגן תמונות שהצוות רשאי להשתמש בהן, לבחור גרסאות ואז לבדוק עקביות. האם המרקם של הבד נשמר? האם הבגד שמופיע בגרסה השנייה זהה לזה שיופיע בחנות? כשמשתמשים בייצוג של אדם מזוהה, הצוות צריך לסכם גם אילו גרסאות ובאילו הקשרים מותר להפיק. ההחלטה אם לבנות סביבה דיגיטלית או להשתמש ברקע מחולל תלויה גם בכמות השליטה שהלקוח צריך בכל זווית ופרט.}

\par}

{\fontsize{12.2}{17.2}\selectfont

השוואת עלויות מתחילה באותה יחידת עבודה. הוצאות על לוקיישן יכולות להצטמצם, ובמקביל להתווסף זמן להכנת קובצי מקור, בדיקות עקביות ותיקון תמונות. ייתכנו גם עלויות רישוי של דמות או חומר ועלויות אישור חוזר. \textenglish{Mango} פירטה את הדרך שבה נבנה הקמפיין שלה, בלי לפרסם תקציב השוואתי. השאלה הראשונה היא כמה עולה לסיים ארבע תמונות ראויות לפרסום, כולל אלה שנפסלו בדרך. \SourceRef{7}

\par}
